\documentclass{ctexart}
\usepackage{Note}
\usepackage{unicode-math}
\setCJKmainfont[
  BoldFont = 思源宋体 Bold,
  ItalicFont = 楷体,
]{宋体}
\begin{document}
\section{全同粒子}
\subsection{双粒子体系}
对单粒子而言, $\iPsi(\vec{r},t)$是空间坐标$\vec{r}$和时间$t$的函数,那么双粒子(以及多粒子)系统的状态则是所有粒子的坐标$\vec{r}_i$以及时间$t$的函数$\iPsi\left(\vec{r}_1,\cdots,\vec{r}_n,t\right)$,其演化方式也遵循薛定谔方程:
\[\i\hbar\dfrac{\p\iPsi}{\p t}=\hat{H}\iPsi\]
\indent 对于双粒子体系而言,哈密顿量
\[\hat{H}=-\dfrac{\hbar^2}{2m_1}\nabla_1^2-\dfrac{\hbar^2}{2m_2}\nabla_2^2+V\left(\vec{r}_1,\vec{r}_2,t\right)\]
这里$\nabla$的下标表示对粒子$1$或粒子$2$的坐标的微分.波函数的统计含义是显然的,即$\left|\iPsi\left(\vec{r}_1,\vec{r}_2,t\right)\right|^2$是在$\vec{r}_1$处发现粒子$1$,在$\vec{r}_2$处发现粒子$2$的概率密度函数;它也同样需要满足归一化条件.\\
\indent 对于不含时的势函数,通过分离变量可以获得一组完备的解:
\[\iPsi\left(\vec{r}_1,\vec{r}_2,t\right)=\psi\left(\vec{r}_1,\vec{r}_2\right)\exp\left(-\dfrac{\i Et}{\hbar}\right)\]
其中空间波函数$\psi$满足定态薛定谔方程:
\[-\dfrac{\hbar^2}{2m_1}\nabla_1^2\psi-\dfrac{\hbar^2}{2m_2}\nabla_2^2\psi+V\psi=E\psi\]
这就是我们主要研究的问题.一般而言,上述方程是困难的,但有两种特殊情况可以把它转化为单粒子问题:
\begin{enumerate}[label=\tbf{\arabic*.},leftmargin=*]
    \item \tbf{无相互作用粒子}\ 假定两个粒子之间没有相互作用,因此势函数可以写作两者势能之和:
    \[V\left(\vec{r}_1,\vec{r}_2\right)=V_1\left(\vec{r}_1\right)+V_2\left(\vec{r}_2\right)\]
    对定态波函数进行变量分离:
    \[\psi\left(\vec{r}_1,\vec{r}_2\right)=\psi_1\left(\vec{r}_1\right)\psi_2\left(\vec{r}_2\right)\]
    这样可以将定态薛定谔方程改写为
    \[-\dfrac{\hbar^2}{2m_i}\nabla_i^2\psi_i\left(\vec{r}_i\right)+V_i\left(\vec{r}_i\right)\psi_i\left(\vec{r}_i\right)=E_i\psi_i\left(\vec{r}_i\right),\quad i=1,2\]
    并且$E=E_1+E_2$.在这种情况下,双粒子波函数是两个单粒子波函数的简单乘积:
    \[\iPsi\left(\vec{r}_1,\vec{r}_2,t\right)=\iPsi_1\left(\vec{r}_1,t\right)\iPsi_2\left(\vec{r}_2,t\right)\]
    另外,这些解的任何线性组合也是成立的解.例如:
    \[\iPsi\left(\vec{r}_1,\vec{r}_2,t\right)=p_1\iPsi_{1,a}\left(\vec{r}_1,t\right)\iPsi_{2,b}\left(\vec{r}_2,t\right)+p_2\iPsi_{1,c}\left(\vec{r}_1,t\right)\iPsi_{2,d}\left(\vec{r}_2,t\right)\]
    这时,测量粒子$1$的能量有$\left|p_1\right|^2$的概率得到$E_a$,此时粒子$2$的能量为$E_b$;也有$\left|p_2\right|^2$的概率得到$E_c$,此时粒子$2$的能量为$E_d$.两个粒子的状态是相互依赖的,我们称这两个粒子是\tbf{纠缠}的.纠缠态不能写作单粒子状态的乘积.
    \item \tbf{中心势场}\ 如果粒子之间的相互作用仅限于二者之间,并且其相互作用的势仅取决于它们的距离,即
    \[V\left(\vec{r}_1,\vec{r}_2\right)\to V\left(\left|\vec{r}_1-\vec{r}_2\right|\right)\]
    那么也可以把系统转化为无相互作用粒子.定义相对位移$\vec{r}=\vec{r}_1-\vec{r}_2$,质心位置$\vec{R}=\dfrac{m_1\vec{r}_1+m_2\vec{r}_2}{m_1+m_2}$和约化质量$\mu=\dfrac{m_1m_2}{m_1+m_2}$.
    于是
    \[\nabla_1=\dfrac{\p}{\p\vec{r}_1}=\dfrac{\p}{\p\vec{r}}\dfrac{\p\vec{r}}{\p\vec{r}_1}+\dfrac{\p}{\p\vec{R}}\dfrac{\p\vec{R}}{\p\vec{r}_1}=\nabla_r+\dfrac{\mu}{m_2}\nabla_R,\quad\nabla_2=-\nabla_r+\dfrac{\mu}{m_1}\nabla_R\]
    这样就有
    \[\dfrac{\hbar^2}{2m_1}\nabla_1^2+\dfrac{\hbar^2}{2m_2}\nabla_2^2=\dfrac{\hbar^2}{2}\left(\dfrac{\nabla_r^2}{m_1}+\dfrac{2\mu\nabla_r\nabla_R}{2m_1m_2}+\dfrac{\mu^2\nabla_R^2}{m_1m_2^2}+\dfrac{\nabla_r^2}{m_2}-\dfrac{2\mu\nabla_r\nabla_R}{2m_1m_2}+\dfrac{\mu^2\nabla_R^2}{m_1^2m_2}\right)=\dfrac{\hbar^2}{2}\left(\dfrac{\nabla_r^2}{\mu}+\dfrac{\nabla_R^2}{m_1+m_2}\right)\]
    于是定态薛定谔方程可以改写为
    \[-\dfrac{\hbar^2}{2\left(m_1+m_2\right)}\nabla_R^2\psi-\dfrac{\hbar^2}{2\mu}\nabla_r^2\psi+V(\vec{r})\psi=E\psi\]
    分离变量,令
    \[\psi\left(\vec{R},\vec{r}\right)=\psi_R(\vec{R})\psi_r(\vec{r})\]
    注意到$\psi_R$满足总质量为$m_1+m_2$,势能为$0$,能量为$E_R$的单粒子薛定谔方程; $\psi_r$满足约化质量为$\mu$,势能为$V(\vec{r})$,能量为$E_r$的单粒子薛定谔方程.这告诉我们质心的运动可以看作一个自由粒子,而两个粒子的相对运动可以看作质量为约化质量的处于中心势场$V(\vec{r})$的单粒子的运动.
\end{enumerate}
\subsubsection{玻色子和费米子}
重新审视我们通过变分法得到的波函数
\[\psi\left(\vec{r}_1,\vec{r}_2\right)=\psi_1\left(\vec{r}_1\right)\psi_2\left(\vec{r}_2\right)\]
这隐含了一个前提,即我们能够把粒子区分开;而实际上我们讨论的粒子是不可分的(即\tbf{全同}的).量子力学通过构造一个波函数来描述不可分的粒子:
\[\psi_{\pm}\left(\vec{r}_1,\vec{r}_2\right)=A\left[\psi_a\left(\vec{r}_1\right)\psi_b\left(\vec{r}_2\right)\pm\psi_b\left(\vec{r}_1\right)\psi_a\left(\vec{r}_2\right)\right]\]
这一理论将允许两种类型的全同粒子:\tbf{玻色子},上式取正号;\tbf{费米子},上式取符号.玻色子是交换对称的,即$\psi_+\left(\vec{r}_1,\vec{r}_2\right)=\psi_+\left(\vec{r}_2,\vec{r}_1\right)$;而费米子是交换反对称的,即$\psi_-\left(\vec{r}_1,\vec{r}_2\right)=-\psi_-\left(\vec{r}_2,\vec{r}_1\right)$.巧合的是,所有自旋为整数的粒子都是玻色子,而所有自旋为半整数的粒子都是费米子(可以用相对论量子力学进行证明,这里不进行过多讨论).\\
\indent 由上面的构造可以得到一个著名的原理:
\begin{theorem}[泡利不相容原理]
    相同的费米子不能占据相同的状态.
\end{theorem}
\noindent 如果两个费米子处于相同的状态,那么$\psi_-\left(\vec{r}_1,\vec{r}_2\right)=A\left[\psi_a\left(\vec{r}_1\right)\psi_a\left(\vec{r}_2\right)-\psi_a\left(\vec{r}_1\right)\psi_a\left(\vec{r}_2\right)\right]=0$,这是没有意义的波函数,因而也不被允许.
\subsubsection{交换力}
粒子的交换对称性将影响它们作用的方式.我们从一个简单的问题来探究前述两种粒子的区别.
\begin{problem}
    给定归一化的一维波函数$\psi_a$, $\psi_b$分别代表$a$, $b$两种状态.计算两个全同的玻色子/费米子的距离平方的期望值$\avg{\left(x_1-x_2\right)^2}$.
\end{problem}
\begin{solution}
首先,容易证明两种情况的波函数分别为
\[\psi_+=\dfrac{1}{\sqrt{2}}\left[\psi_a\left(x_1\right)\psi_b\left(x_2\right)+\psi_b\left(x_1\right)\psi_a\left(x_2\right)\right]\]
\[\psi_-=\dfrac{1}{\sqrt{2}}\left[\psi_a\left(x_1\right)\psi_b\left(x_2\right)-\psi_b\left(x_1\right)\psi_a\left(x_2\right)\right]\]
因此
\[\begin{aligned}
    \avg{x_1^2}
\end{aligned}\]
\end{solution}
\end{document}